\section{Final checklist}

This chapter introduced continuous random variables as probability models in
which probability is spread over intervals rather than concentrated at individual
points. The following checklist summarizes the ideas that should be understood
before moving on.

\subsection*{Conceptual understanding}

You should be able to explain the following points in words:
\begin{itemize}
    \item why a continuous random variable can satisfy \(P(X=x)=0\) for every
    individual value \(x\), while still taking values in a whole interval;
    \item why probabilities of intervals, not probabilities of individual
    points, are the fundamental objects in continuous models;
    \item why a density value \(f_X(x)\) is not usually a probability;
    \item why probability is represented by area under a density curve;
    \item why changing endpoints usually does not change probabilities in
    continuous density models;
    \item how the continuous case parallels the discrete case, and where the
    analogy must be treated carefully.
\end{itemize}

\subsection*{Definitions and formulas}

You should know the meaning of the following formulas:
\[
P(a\leq X\leq b)=\int_a^b f_X(x)\,dx,
\]
\[
\int_{-\infty}^{\infty}f_X(x)\,dx=1,
\]
\[
F_X(x)=P(X\leq x)=\int_{-\infty}^x f_X(t)\,dt,
\]
\[
P(a\leq X\leq b)=F_X(b)-F_X(a),
\]
\[
\E[X]=\int_{-\infty}^{\infty}x f_X(x)\,dx,
\]
\[
\Var(X)=\E[(X-\E[X])^2]=\E[X^2]-(\E[X])^2.
\]

You should also remember that these formulas require the relevant integrals to
exist. In particular, expectation and variance are not guaranteed to be finite
for every possible distribution.

\subsection*{Typical tasks}

After studying this chapter, you should be able to:
\begin{itemize}
    \item check whether a proposed function is a valid density;
    \item compute interval probabilities by integrating a density;
    \item compute interval probabilities by using a CDF;
    \item derive a CDF from a density in simple examples;
    \item recover a density from a differentiable CDF;
    \item compute expectation and variance for simple continuous distributions;
    \item find quantiles by solving \(F_X(q_p)=p\) when the CDF is strictly
    increasing;
    \item interpret the support, shape, and parameters of a continuous model.
\end{itemize}

\subsection*{Common mistakes}

The most common mistakes in this topic are not algebraic; they are conceptual.
Avoid the following errors:
\begin{itemize}
    \item treating \(f_X(x)\) as if it were \(P(X=x)\);
    \item thinking that a density cannot be larger than \(1\);
    \item forgetting to check that the total area under a density is \(1\);
    \item using sums where integrals are needed;
    \item assuming that probability zero always means impossibility;
    \item confusing the expectation with the most likely value or with the
    median;
    \item choosing a classical distribution by name without checking whether its
    support and modelling assumptions match the situation.
\end{itemize}

\subsection*{Questions for self-explanation}

A good test of understanding is whether you can answer the following questions
without immediately searching for a formula:
\begin{enumerate}
    \item If \(P(X=x)=0\) for every \(x\), how can \(X\) still have total
    probability \(1\)?
    \item What is the difference between the height of a density and the area
    under a density?
    \item Why does \(P(a<X<b)=P(a\leq X\leq b)\) for the density models studied
    here?
    \item How is a CDF built from a density?
    \item How is an interval probability read from a CDF?
    \item Why is expectation an average even though no point has positive
    probability?
    \item What does variance measure geometrically or intuitively?
    \item What modelling situation might suggest a uniform, exponential, normal,
    gamma, or beta distribution?
\end{enumerate}

\begin{summarybox}
The essential message of the chapter is this: in continuous probability, point
probabilities disappear, but probability itself does not disappear. It is
accumulated over intervals as area under a density, and the familiar ideas of
distribution, expectation, variance, and quantiles are rebuilt using integrals.
\end{summarybox}
